% ================================================================
%  §2  Linear Algebra and Quadratic Forms
% ================================================================
\section{Linear Algebra and Quadratic Forms}\label{sec:spectral}

The nicest objectives are linear and quadratic --- and even these require nontrivial algebra once the dimension grows.
This section climbs a ladder of complexity: linear functions (trivial), univariate quadratics (trivial via a trick), separable multivariate quadratics (still easy), and finally general quadratics (where variables interact and the real work begins).
At the top of the ladder sits spectral theory, which reveals that eigenvalues control everything about a quadratic: curvature, level-set geometry, solvability, and computational cost.


% ────────────────────────────────────────────────────────────────
\subsection{Scalar Products, Norms, and Distances}\label{ssec:inner-product}

Before analyzing any function on $\R^n$, we need the geometric primitives: how to measure angles, lengths, and proximity.

\begin{defn}[Scalar product]\label{def:scalar-product}
	The \textit{(Euclidean) scalar product} of $x, z \in \R^n$ is
	$\inner{x, z} = \sum_{i=1}^{n} x_i z_i$.
	It satisfies four defining properties:
	symmetry ($\inner{x, z} = \inner{z, x}$),
	positive definiteness ($\inner{x, x} \geq 0$ with equality iff $x = 0$),
	homogeneity ($\inner{\alpha x, z} = \alpha \inner{x, z}$),
	and linearity ($\inner{x + w, z} = \inner{x, z} + \inner{w, z}$).
	Other scalar products exist and make sense in other spaces (matrices, integrable functions, random variables); cf.\ the kernel trick in SVMs (\Cref{sec:constrained-theory}).
\end{defn}

The scalar product encodes direction: $\inner{x, z} = \norm{x} \cdot \norm{z} \cdot \cos(\theta)$, where $\theta$ is the angle between $x$ and $z$.
Positive means ``same direction,'' zero means orthogonal ($x \perp z$), negative means ``opposite direction.''
The \textit{Cauchy--Schwarz inequality} $\abs{\inner{x, z}} \leq \norm{x} \norm{z}$ is an immediate consequence.

\begin{defn}[Norm and distance]\label{def:norm-distance}
	The \textit{(Euclidean) norm} of $x \in \R^n$ is $\norm{x} = \sqrt{\inner{x, x}} = \sqrt{x_1^2 + \cdots + x_n^2}$.
	It satisfies: $\norm{x} \geq 0$ with equality iff $x = 0$; $\norm{\alpha x} = \abs{\alpha} \norm{x}$; and the \textit{triangle inequality} $\norm{x + z} \leq \norm{x} + \norm{z}$.
	Two related identities: $\norm{x + z}^2 = \norm{x}^2 + \norm{z}^2 + 2\inner{x, z}$ and the \textit{parallelogram law} $2\norm{x}^2 + 2\norm{z}^2 = \norm{x + z}^2 + \norm{x - z}^2$.
	The \textit{distance} between $x$ and $z$ is $d(x, z) = \norm{x - z}$.
\end{defn}

\begin{defn}[Ball]\label{def:ball}
	The closed ball centered at $x \in \R^n$ with radius $r > 0$ is
	$\mathcal{B}(x, r) = \{z \in \R^n : \norm{z - x} \leq r\}$.
	Balls replace intervals as the basic notion of ``neighborhood'' in $\R^n$.
	The norm determines the topology of the space --- what is ``close'' to what --- and this will be crucial when we discuss convergence and continuity.
\end{defn}


The Euclidean norm is just one member of a large family.
The \textit{$p$-norm} is $\norm{x}_p = (\sum_{i=1}^{n} |x_i|^p)^{1/p}$ for $p \geq 1$.
Important special cases: $\norm{x}_2$ (Euclidean), $\norm{x}_1 = \sum_i |x_i|$ (used in Lasso regularization, \Cref{ssec:regularization}), $\norm{x}_\infty = \max_i |x_i|$ (Chebyshev), and $\norm{x}_0 = \#\{i : |x_i| > 0\}$ (counting ``norm'' --- not actually a norm, since it violates homogeneity).
Different $p$ give different ball shapes: $\mathcal{B}_1$ is a diamond, $\mathcal{B}_2$ a sphere, $\mathcal{B}_\infty$ a hypercube.
In finite dimension all norms are \textit{equivalent}: for any two norms $\norm{\cdot}$ and $\norm{\cdot}'$, there exist $0 < \alpha \leq \beta$ such that $\alpha\norm{x} \leq \norm{x}' \leq \beta\norm{x}$ for all $x$.
This means convergence in one norm implies convergence in every other --- the topology of $\R^n$ doesn't depend on the choice of norm.


% ────────────────────────────────────────────────────────────────
\subsection{Matrices}\label{ssec:matrices}

A matrix $A \in \R^{m \times n}$ has $m$ rows $A_i \in \R^n$ and $n$ columns $A^j \in \R^m$.
The \textit{transpose} $A^T \in \R^{n \times m}$ exchanges rows with columns: $(A^T)_{ij} = A_{ji}$.
A square matrix $A \in \R^{n \times n}$ is \textit{symmetric} if $A = A^T$, equivalently $A_{ij} = A_{ji}$ for all $i, j$.
A scalar $a \in \R$ can be viewed as a $1 \times 1$ symmetric matrix.

$\R^n$ is a finite-dimensional vector space of dimension $n$.
The \textit{canonical basis} $\{u^1, \dots, u^n\}$ consists of the $n$ unit vectors: $u^i_j = 1$ if $j = i$ and $0$ otherwise.
Every $x \in \R^n$ is a linear combination $x = \sum_i x_i u^i$.

The matrix--vector product $Ax \in \R^m$ (for $x \in \R^n$) has $i$-th component $\inner{A_i, x}$; the product $x^T A$ (for $x \in \R^m$) gives a row vector with $j$-th component $\inner{A^j, x}$.
The scalar product can be written $\inner{x, z} = x^T z = z^T x$.
The matrix--matrix product $AB$ (for $A \in \R^{m \times k}$, $B \in \R^{k \times n}$) has entry $(AB)_{ij} = \inner{A_i, B^j}$ and costs $\mathcal{O}(nmk)$ --- expensive unless $A$ or $B$ is sparse.

The \textit{spectral norm} (or \textit{operator norm}) of $A$ is
$\norm{A} = \max_{\norm{x}=1} \norm{Ax} = \sigma_{\max}(A) = \sqrt{\lambda_{\max}(A^T A)}$,
where $\sigma_{\max}$ is the largest singular value.
For symmetric $A$, $\norm{A} = \max_i |\lambda_i|$.
The spectral norm is submultiplicative ($\norm{AB} \leq \norm{A}\norm{B}$) and controls how much a matrix can stretch a vector.
The \textit{spectral radius} $\rho(A) = \max_i |\lambda_i(A)|$ satisfies $\rho(A) \leq \norm{A}$, with equality when $A$ is symmetric.


% ────────────────────────────────────────────────────────────────
\subsection{Linear Functions}\label{ssec:linear-functions}

The simplest possible objective is linear.
We start univariate, then generalize.

\paragraph{Univariate case}

A \textit{linear univariate function} is $f(x) = bx$ for a fixed $b \in \R$.
The parameter $b$ is the \textit{slope}: if $b > 0$, $f$ is increasing; if $b < 0$, decreasing; if $b = 0$, constant.
Unconstrained optimization is trivial: $\min f = -\infty$ and $\max f = +\infty$ unless $b = 0$, in which case both equal $0$.

Box-constrained optimization $\min\{bx : x \in [x_-, x_+]\}$ is equally immediate.
When $b > 0$, the minimum sits at $x_-$ and the maximum at $x_+$; when $b < 0$ the roles reverse; when $b = 0$ every point is optimal.
This ``works'' even if $x_- = -\infty$ or $x_+ = +\infty$, since $b \cdot (\pm \infty) = \pm \infty$ by convention.
A closed-form $\mathcal{O}(1)$ solution --- don't get used to it.

\paragraph{Multivariate case}

A \textit{linear multivariate function} is $f(x) = \inner{b, x} = \sum_{i=1}^{n} b_i x_i$ for a fixed $b \in \R^n$.
Linearity means $f(\gamma x) = \gamma f(x)$ and $f(x + z) = f(x) + f(z)$.
The graph of $f$ is a hyperplane in $\R^{n+1}$; level sets are parallel hyperplanes in $\R^n$, all perpendicular to $b$.

For box-constrained optimization on a hyperrectangle $X = \{x : x_- \leq x \leq x_+\}$, the problem separates into $n$ independent univariate problems (one per coordinate), each solvable in $\mathcal{O}(1)$.
Total cost: $\mathcal{O}(n)$.
This is the last time optimization will be this easy.


% ────────────────────────────────────────────────────────────────
\subsection{Quadratic Univariate Functions}\label{ssec:quad-univariate}

The next step up: $f(x) = ax^2 + bx$ for fixed $a, b \in \R$.
A homogeneous quadratic ($b = 0$) plus a linear ($a = 0$) term.

\paragraph{Homogeneous case ($b = 0$)}

$f(x) = ax^2$ is symmetric around the origin.
If $a > 0$, the parabola opens upward: $f$ decreases for $x < 0$, increases for $x > 0$, and the minimum is at $x = 0$.
If $a < 0$, it opens downward: the maximum is at $x = 0$, and $f$ is unbounded below.
If $a = 0$, $f$ is constant.

Box-constrained optimization on $[x_-, x_+]$ requires a case analysis.
When $a > 0$: if $0 \in [x_-, x_+]$ then $\argmin = 0$ and $\argmax = \argmax\{f(x_-), f(x_+)\}$; if $x_- > 0$ then $\argmin = x_-$; if $x_+ < 0$ then $\argmin = x_+$.
Still $\mathcal{O}(1)$.

\paragraph{Non-homogeneous case ($a \neq 0$, $b \neq 0$)}

The linear term $bx$ breaks the symmetry around the origin.
A powerful trick restores it: define $\bar{x} = -b / (2a)$ and substitute $z = x - \bar{x}$.
Expanding $f(x) = a(z + \bar{x})^2 + b(z + \bar{x})$ and using $2a\bar{x} + b = 0$, the cross-terms cancel:
\begin{equation}\label{eq:completing-square-univariate}
	f(x) = az^2 + f(\bar{x}).
\end{equation}
The problem reduces to a homogeneous quadratic in $z$ plus a constant offset.
The point $\bar{x}$ is the \textit{center} of the function --- the vertex of the parabola.
If $a > 0$, the unique minimum is at $\bar{x}$; if $a < 0$, the unique maximum.
Still $\mathcal{O}(1)$.

\begin{rem}\label{rem:completing-square}
	This ``completing the square'' trick will reappear in the multivariate setting (\Cref{ssec:quad-nonhom}).
	The idea is always the same: shift the origin to the center, where the function becomes homogeneous and therefore easier to analyze.
\end{rem}


% ────────────────────────────────────────────────────────────────
\subsection{Separable Quadratic Functions}\label{ssec:separable-quad}

A \textit{separable (non-homogeneous) quadratic} is
\begin{equation}\label{eq:separable-quad}
	f(x) = \sum_{i=1}^{n} \bigl(a_i x_i^2 + b_i x_i\bigr),
	\qquad (a, b) \in \R^{2n}.
\end{equation}
This is a sum of $n$ univariate quadratics --- one per coordinate --- with no cross-terms.
Each can be minimized independently by completing the square: the center is $\bar{x}_i = -b_i / (2a_i)$ when $a_i \neq 0$.
Total cost: $\mathcal{O}(n)$.

An important special case is $f(x) = \norm{x}^2 = \sum_i x_i^2$ (all $a_i = 1$, $b_i = 0$), whose level sets are circles (when $n = 2$) or spheres.
When $a_1 \neq a_2$, the level sets become axis-aligned ellipses: the larger $a_i / a_j$, the more elongated.
Non-zero $b_i$ simply shifts the center from the origin to $\bar{x}$.

This is the last time we get closed formulas.
The separable structure means variables don't interact --- each $x_i$ can be chosen without knowing the others.
In a general quadratic, this independence breaks.


% ────────────────────────────────────────────────────────────────
\subsection{Tomographies}\label{ssec:tomographies}

Before tackling the general multivariate case, we need a tool to visualize and analyze functions of $n$ variables.
The graph $\Gr(f) \subset \R^{n+1}$ is unpicturable for $n > 2$.
Level sets help but live in $\R^n$, still hard for $n > 3$.
The solution: ``slice'' the function along a chosen direction and study the resulting one-dimensional function.

\begin{defn}[Tomography]\label{def:tomography}
	Given $f \colon \R^n \to \R$, an origin $x \in \R^n$, and a direction $d \in \R^n$, the \textit{tomography} of $f$ along $d$ from $x$ is
	\begin{equation}\label{eq:tomography}
		\varphi_{x,d}(\alpha) = f(x + \alpha d) \colon \R \to \R.
	\end{equation}
	When $x$ and $d$ are clear from context, we write simply $\varphi(\alpha)$.
\end{defn}

The norm of $d$ only changes the horizontal scale: $\varphi_{x, \beta d}(\alpha) = \varphi_{x, d}(\beta \alpha)$, so it is often convenient to use a normalized direction $\norm{d} = 1$.
Restricting to a coordinate direction $u^i$ (the $i$-th canonical basis vector) gives the partial restriction $f_x^i(\alpha) = f(x_1, \dots, x_{i-1}, \alpha, x_{i+1}, \dots, x_n)$.

\paragraph{Tomography of a linear function}

For $f(x) = \inner{b, x}$ with origin $x = 0$ and $\norm{d} = 1$:
\begin{equation}\label{eq:tomography-linear}
	\varphi(\alpha) = \alpha \inner{b, d} = \alpha \norm{b} \cos(\theta),
\end{equation}
where $\theta$ is the angle between $b$ and $d$.
The tomography is a linear function of $\alpha$.
It is steepest when $d$ is collinear with $b$ (either same direction or opposite); flat when $d \perp b$.
The direction $-b / \norm{b}$ gives the steepest decrease --- a foreshadowing of the gradient method.

Line search methods exploit tomographies directly: pick a descent direction $d$, then minimize $\varphi_{x,d}$ over $\alpha \geq 0$.
We will see this idea at work in every gradient-based algorithm.


% ────────────────────────────────────────────────────────────────
\subsection{The General Quadratic Function}\label{ssec:general-quad}

Now the variables interact.
The general (homogeneous) quadratic on $\R^n$ is
\begin{equation}\label{eq:homogeneous-quad}
	f(x) = \frac{1}{2}\, x^T Q x
	= \frac{1}{2} \biggl[\sum_{i=1}^{n} Q_{ii} x_i^2
		+ \sum_{\substack{i,j=1 \\ j \neq i}}^{n} Q_{ij} x_i x_j\biggr],
\end{equation}
for a fixed $Q \in \R^{n \times n}$.
The off-diagonal terms $Q_{ij} x_i x_j$ couple the variables: changing $x_i$ affects the optimal choice for $x_j$.
This coupling is what makes general quadratics fundamentally harder than separable ones.

An important observation: even if $Q$ is not symmetric, we can always replace it with $(Q + Q^T)/2$ without changing $f$.
Indeed, $x^T Q x = x^T [(Q + Q^T)/2]\, x$.
So there is no loss of generality in assuming $Q$ symmetric from the start.

The tomography along a direction $d$ from the origin gives
\begin{equation}\label{eq:tomography-quad}
	\varphi(\alpha) = f(\alpha d) = \frac{1}{2}\, \alpha^2 (d^T Q d).
\end{equation}
This is a univariate parabola whose sign and steepness depend entirely on $d^T Q d$.
Different directions yield different curvatures.
Understanding this dependence requires knowing the spectrum of $Q$.


% ────────────────────────────────────────────────────────────────
\subsection{Eigenvalues and Spectral Decomposition}\label{ssec:spectral}

\begin{defn}[Eigenvalues and eigenvectors]\label{def:eigen}
	Given a matrix $Q \in \R^{n \times n}$, a scalar $\lambda \in \R$ is an \textit{eigenvalue} and a nonzero vector $v \in \R^n$ is the corresponding \textit{eigenvector} if $Qv = \lambda v$.
	An eigenvector is a direction along which $Q$ acts as a simple rescaling.
\end{defn}

Eigenvalues can be complex in general, and eigenvectors need not be orthogonal.
Symmetric matrices are the exception --- and the one that matters for optimization.

\begin{thm}[Spectral theorem for symmetric matrices]\label{thm:sym-eigen}
	If $Q \in \R^{n \times n}$ is symmetric, then $Q$ has $n$ real eigenvalues (counted with multiplicity) and $n$ eigenvectors that can be chosen orthonormal.
\end{thm}

We denote these orthonormal eigenvectors $H_1, \dots, H_n \in \R^n$ with corresponding eigenvalues $\lambda_1 \geq \lambda_2 \geq \cdots \geq \lambda_n$ (ordered from largest to smallest), so that $Q H_i = \lambda_i H_i$ for all $i$.

\begin{defn}[Eigenvector matrix and spectral decomposition]\label{def:spectral-decomp}
	Collecting the eigenvectors into $H = [H_1 \mid \cdots \mid H_n] \in \R^{n \times n}$, orthonormality gives $H^T H = I$ (so $H$ is orthogonal).
	The \textit{spectral decomposition} is
	\begin{equation}\label{eq:spectral-decomp}
		Q = H \Lambda H^T = \sum_{i=1}^{n} \lambda_i H_i H_i^T,
		\qquad \Lambda = \diag(\lambda_1, \dots, \lambda_n).
	\end{equation}
	Each term $\lambda_i H_i H_i^T$ is a rank-one matrix.
	The decomposition expresses $Q$ as a weighted sum of these terms, one per eigenvalue.
\end{defn}

Why does this matter?
The spectral decomposition turns matrix operations into scalar ones.
If we change coordinates via $y = H^T x$, the quadratic $x^T Q x$ becomes $\sum_i \lambda_i y_i^2$ --- a sum of independent squared terms, each scaled by an eigenvalue.
Large eigenvalues stretch, small ones compress.
The ratio between the largest and smallest controls how ``distorted'' the quadratic is, and (as we will see) how fast gradient methods converge.

\begin{thm}[Variational characterization of eigenvalues]\label{thm:var-eigen}
	For a symmetric $Q$, the extreme eigenvalues satisfy
	\begin{equation}\label{eq:rayleigh-quotient}
		\lambda_1 = \max\left\{\frac{d^T Q d}{d^T d} : d \neq 0\right\},
		\qquad
		\lambda_n = \min\left\{\frac{d^T Q d}{d^T d} : d \neq 0\right\}.
	\end{equation}
	The ratio $d^T Q d / d^T d$ is the \textit{Rayleigh quotient}.
	The maximum is attained at $d = H_1$, the minimum at $d = H_n$.
\end{thm}

\paragraph{Determinant and inverse}

\begin{defn}[Determinant]\label{def:determinant}
	The determinant of $Q$ equals the product of its eigenvalues:
	$\det(Q) = \prod_{i=1}^{n} \lambda_i$.
	If $\det(Q) \neq 0$, then $Q$ is nonsingular, every eigenvalue is nonzero, and $Q^{-1}$ exists.
	The eigenvectors of $Q^{-1}$ are the same as those of $Q$, with reciprocal eigenvalues: $Qv = \lambda v$ implies $Q^{-1}v = (1/\lambda)\, v$.
	It follows that $\det(Q^{-1}) = \det(Q)^{-1}$.
\end{defn}

Computing eigenvalues by finding roots of the characteristic polynomial $\det(Q - \lambda I) = 0$ is practical only for $n = 2$ (quadratic formula) or $n = 3$ (cubic formula).
For $n = 2$: $\det(Q) = Q_{11}Q_{22} - Q_{12}^2$ (using symmetry).
For larger $n$, iterative algorithms (QR, divide-and-conquer) are used instead \cite{nocedal2006}.
Diagonal and triangular matrices are the exception: their eigenvalues are the diagonal entries.


% ────────────────────────────────────────────────────────────────
\subsection{Matrix Definiteness}\label{ssec:definiteness}

The sign pattern of the eigenvalues determines the shape of the associated quadratic --- and therefore the nature of its critical points.

\begin{defn}[Matrix definiteness]\label{def:definiteness}
	A symmetric matrix $Q$ is:
	\begin{itemize}
		\item \textit{Positive definite} ($Q \succ 0$): all $\lambda_i > 0$.
			The quadratic $\frac{1}{2} x^T Q x$ has a unique global minimum at the origin.
		\item \textit{Negative definite} ($Q \prec 0$): all $\lambda_i < 0$.
			Unique global maximum at the origin.
		\item \textit{Positive semidefinite} ($Q \succeq 0$): all $\lambda_i \geq 0$.
			The quadratic is convex, but the minimum need not be unique --- flat directions correspond to zero eigenvalues.
		\item \textit{Negative semidefinite} ($Q \preceq 0$): all $\lambda_i \leq 0$.
			Concave, with a possibly non-unique maximum.
		\item \textit{Indefinite}: eigenvalues of both signs.
			The origin is a saddle point.
	\end{itemize}
\end{defn}


% ────────────────────────────────────────────────────────────────
\subsection{Tomographies of Quadratic Functions}\label{ssec:quad-tomographies}

The spectral decomposition makes tomographies along eigenvectors completely transparent.

\begin{equation}\label{eq:tomography-eigenvector}
	\varphi_{0, H_i}(\alpha) = f(\alpha H_i) = \frac{1}{2}\, \alpha^2 \lambda_i.
\end{equation}
Along eigenvector $H_i$, the quadratic reduces to a univariate parabola scaled by $\lambda_i$.
The steepest upward curvature is along $H_1$ ($\lambda_1$, the largest eigenvalue); the steepest downward curvature (or least upward) is along $H_n$ ($\lambda_n$, the smallest).
Intermediate directions give intermediate curvatures, always in $[\lambda_n,\lambda_1]$.

\begin{exm}[Tomographies under different definiteness]\label{exm:quad-tomographies}
	Consider $n = 2$ with $H = \frac{\sqrt{2}}{2}\bigl[\begin{smallmatrix} -1 & 1 \\ 1 & 1 \end{smallmatrix}\bigr]$ (fixed eigenvectors, rotated $45^\circ$).
	\begin{itemize}
		\item $Q \succ 0$ with $\lambda = (8, 4)$: $d^T Q d > 0$ for all $d \neq 0$.
			Every tomography is an upward parabola.
			Steepest along $H_1$ ($\lambda_1 = 8$); least steep along $H_2$ ($\lambda_2 = 4$).
		\item $Q \succeq 0$ with $\lambda = (8, 0)$: $d^T Q d \geq 0$, but the tomography along $H_2$ is completely flat.
			The quadratic has a ``valley'' in the $H_2$-direction.
		\item $Q$ indefinite with $\lambda = (8, -2)$: $d^T Q d$ can be positive or negative depending on $d$.
			Along $H_1$ the parabola opens upward; along $H_2$ it opens downward.
			The origin is a saddle point.
		\item $Q \prec 0$ with $\lambda = (-4, -8)$: $d^T Q d < 0$ for all $d \neq 0$.
			Every tomography is a downward parabola.
			Steepest negative curvature along $H_2$ ($\lambda_2 = -8$).
	\end{itemize}
\end{exm}


% ────────────────────────────────────────────────────────────────
\subsection{Level Sets and Geometry}\label{ssec:level-sets}

The spectral decomposition also governs the shape of level sets.
For a homogeneous quadratic $f(x) = \frac{1}{2} x^T Q x$, all level sets are centered at the origin (since $f$ is symmetric: $f(x) = f(-x)$).

The level set $L(f, v)$ intersects each eigendirection $H_i$ at distance $\sqrt{2v / \lambda_i}$ from the origin (provided $\lambda_i > 0$ and $v > 0$).
So the eigenvectors are the axes of the level-set ellipsoid, and the axis lengths are $\sqrt{2v / \lambda_i}$.
Smaller $\lambda_i$ means a longer axis: the function is flatter in that direction, so it takes more distance to reach a given function value.

Four geometric regimes correspond to the four definiteness classes:
\begin{itemize}
	\item All $\lambda_i > 0$ (positive definite): level sets are ellipsoids.
		When $\lambda_1 = \lambda_n$ (identity), they are spheres.
	\item Some $\lambda_i = 0$ (positive semidefinite): level sets are ``degenerate'' ellipsoids --- cylinders extending to infinity along the null-space directions.
	\item Eigenvalues of both signs (indefinite): level sets are hyperboloids.
		The positive-eigenvalue directions give the ``bowl'' part; the negative-eigenvalue directions give the ``saddle'' part.
	\item All $\lambda_i < 0$ (negative definite): level sets are ellipsoids again (the function is a downward paraboloid).
\end{itemize}


% ────────────────────────────────────────────────────────────────
\subsection{Optimizing Quadratic Functions}\label{ssec:opt-quad}

Definiteness of $Q$ and the presence of a linear term jointly determine whether --- and how --- a quadratic can be optimized.

\subsubsection{Homogeneous Multivariate Quadratics}\label{sssec:opt-quad-hom}

For $f(x) = \frac{1}{2} x^T Q x$ (no linear term), optimality depends entirely on the eigenvalue signs:
\begin{itemize}
	\item $Q \succeq 0$ and $Q \preceq 0$ simultaneously (i.e., $Q = 0$): $f \equiv 0$, $\min = \max = 0$.
	\item $Q \succeq 0$: $\min = 0$ at $x = 0$, but $\max = +\infty$.
	\item $Q \preceq 0$: $\max = 0$ at $x = 0$, but $\min = -\infty$.
	\item $Q$ indefinite: both $\min = -\infty$ and $\max = +\infty$.
\end{itemize}

What about box-constrained optimization $\min\{f(x) : x \in [x_-, x_+]\}$?
Unlike the univariate case, this is \emph{not} easy in general.
Optimizing a quadratic over a box is $\mathcal{NP}$-hard for indefinite $Q$ \cite{deklerk2008}.
Even when $Q \succeq 0$, the problem is polynomial but nontrivial --- the algorithms we develop in \Cref{sec:quadratic-methods} are needed.
Minimization and maximization behave very differently here.

\subsubsection{Non-homogeneous, Nonsingular Case}\label{ssec:quad-nonhom}

The full quadratic is
\begin{equation}\label{eq:full-quad}
	f(x) = \frac{1}{2}\, x^T Q x + \inner{q, x} + c,
	\qquad Q \in \R^{n \times n},\; q \in \R^n,\; c \in \R.
\end{equation}

When $Q$ is nonsingular ($\det(Q) \neq 0$, equivalently $\lambda_i \neq 0$ for all $i$), we can repeat the univariate trick.
Setting $\nabla f(x) = Qx + q = 0$ yields the critical point $\bar{x} = -Q^{-1} q$.
Think of $\bar{x}$ as the center of the function.
Substituting $z = x - \bar{x}$ and using $Q\bar{x} = -q$, all linear terms cancel:
\begin{equation}\label{eq:completing-square-multivariate}
	f(x) = \frac{1}{2}\, z^T Q z + f(\bar{x}).
\end{equation}
The problem reduces to a homogeneous quadratic in $z$ plus a constant offset.
If $Q \succ 0$, the unique minimum sits at $z = 0$, i.e., $x = \bar{x}$.
If $Q$ is indefinite, $\bar{x}$ is a saddle point.
This is the exact multivariate analogue of the completing-the-square trick.

\subsubsection{Non-homogeneous, Singular Case}\label{sssec:opt-quad-singular}

When $Q$ is singular, $Q^{-1}$ doesn't exist and the trick above fails.
The analysis uses the kernel--image decomposition of $Q$.

Let $I^0 = \{i : \lambda_i = 0\}$ and $I^+ = \{i : \lambda_i \neq 0\}$, with $k = |I^0| > 0$.
The kernel $\ker(Q) = \{v : Qv = 0\}$ is the span of $\{H_i\}_{i \in I^0}$, and the image $\im(Q)$ is the span of $\{H_i\}_{i \in I^+}$.
Since the eigenvectors form an orthonormal basis, $\R^n = \im(Q) \oplus \ker(Q)$ with $\im(Q) \perp \ker(Q)$.

Decompose $q = q^+ + q^0$ with $q^+ \in \im(Q)$ and $q^0 \in \ker(Q)$.
The component $q^+$ lives in the range of $Q$, so we can absorb it via a partial change of variables (finding $\bar{x}$ with $Q\bar{x} = -q^+$).
With $z = x - \bar{x}$:
\begin{equation}\label{eq:singular-quad-reduced}
	f(x) = \frac{1}{2}\, z^T Q z + \inner{q^0, z} + f(\bar{x}).
\end{equation}
The quadratic term $\frac{1}{2} z^T Q z$ kills the kernel directions (because $Qv = 0$ for $v \in \ker(Q)$), but the linear term $\inner{q^0, z}$ may not vanish along those directions.
What matters is $q^0$:

\begin{itemize}
	\item If $q^0 = 0$ (equivalently $q \in \im(Q)$, or $Q\bar{x} = -q$ has a solution): a minimum exists (assuming $Q \succeq 0$), but it is not unique.
		The entire affine subspace $\bar{x} + \ker(Q)$ consists of minimizers --- the function is flat along null-space directions, and all these points share the same objective value $f(\bar{x})$.
	\item If $q^0 \neq 0$ (equivalently $q \notin \im(Q)$): no minimum exists.
		Moving along the null-space direction $q^0$, the quadratic term stays constant but the linear term $\inner{q^0, z}$ keeps decreasing.
		The function is unbounded below.
\end{itemize}

In either case, the key computational task is the same: solve $Q\bar{x} = -q$ (or prove it has no solution).
This is a linear system --- $\mathcal{O}(n^3)$ at worst, doable for moderate $n$, but prohibitive for $n \approx 10^9$.
Iterative methods that avoid forming $Q^{-1}$ become essential, and they are the subject of \Cref{sec:quadratic-methods}.